% !TEX program = xelatex
% ============================================================
%  شرائح محاضرة — Lecture Slides Template
%  قالب شرائح لمحاضرة جامعية مع أهداف تعليمية وأمثلة محلولة
%  Compile with: xelatex main.tex (run twice)
% ============================================================
%  Features:
%    - 16:9 widescreen, clean default beamer theme with accent color
%    - Arabic RTL via polyglossia + Amiri font
%    - Section dividers via \AtBeginSection[]{\frame{\sectionpage}}
%    - Learning objectives, definitions (blocks), worked example
%    - Summary and further reading slides
%    - English terms wrapped with \textEN{...}
% ============================================================

\documentclass[aspectratio=169,11pt]{beamer}

% ---- Arabic / RTL language support ----
\usepackage{polyglossia}
\setdefaultlanguage[calendar=gregorian,hijricorrection=1]{arabic}
\setotherlanguage[variant=british]{english}

% ---- Fonts (beamer uses sans by default, so set sans to Amiri too) ----
\usepackage[no-math]{fontspec}
\setmainfont[Scale=1]{NotoKufiArabic-Regular.ttf}
\setsansfont[Scale=1]{NotoKufiArabic-Regular.ttf}
\newfontfamily\arabicfont[Script=Arabic,Scale=0.95]{NotoKufiArabic-Regular.ttf}
\newfontfamily\arabicfontsf[Script=Arabic,Scale=0.95]{NotoKufiArabic-Regular.ttf}
\newfontfamily\englishfont[Scale=0.95]{TeX Gyre Termes}
\let\textEN\textenglish

% ---- Core packages ----
\usepackage{graphicx}
\usepackage{booktabs}
\usepackage{tabularx}
\usepackage{tikz}
\usepackage{amsmath}
\usepackage{amssymb}

% ---- Theme & appearance ----
\usetheme{default}
\setbeamertemplate{navigation symbols}{}
\setbeamersize{text margin left=10mm, text margin right=10mm}

% ---- Accent color ----
\definecolor{accent}{HTML}{1B5E20}
\definecolor{accentlight}{HTML}{A5D6A7}
\setbeamercolor{frametitle}{fg=accent}
\setbeamercolor{title}{fg=accent}
\setbeamercolor{alerted text}{fg=accent}
\setbeamercolor{structure}{fg=accent}
\setbeamercolor{block title}{fg=white,bg=accent}
\setbeamercolor{block body}{bg=accentlight!20}
\setbeamertemplate{frametitle}[default][center]

% ---- Section divider frames ----
\AtBeginSection[]{
  \begin{frame}[plain]
    \centering
    \vspace{2cm}
    {\Huge\bfseries\color{accent}\insertsectionhead}\\[0.5cm]
    {\large\color{gray!70} القسم \thesection}
  \end{frame}
}

% ---- Metadata ----
% TODO: استبدل ببيانات المقرر والمحاضرة الفعلية
\title{عنوان المقرر}
\subtitle{المحاضرة رقم 1: عنوان المحاضرة}
\author{اسم المدرّس}
\institute{الجامعة}
\date{الفصل الدراسي، سنة}

\begin{document}

% ============================================================
% TITLE SLIDE
% ============================================================
\begin{frame}[plain]
\titlepage
\end{frame}

% ============================================================
% LEARNING OBJECTIVES
% ============================================================
\begin{frame}{أهداف التعلم}
% TODO: حدد ما يجب على الطالب إدراكه بعد المحاضرة
بنهاية هذه المحاضرة، يكون الطالب قادراً على:
\begin{itemize}
  \item شرح المفهوم الأساسي \textEN{(core concept)}.
  \item تطبيق القاعدة في مسائل بسيطة.
  \item تمييز الحالات التي تنطبق عليها القاعدة.
  \item ربط المفهوم بمواضيع لاحقة في المقرر.
\end{itemize}
\end{frame}

% ============================================================
% OUTLINE
% ============================================================
\begin{frame}{المحتويات}
\tableofcontents
\end{frame}

% ============================================================
% SECTION 1 — CONCEPT
% ============================================================
\section{المفاهيم الأساسية}

\begin{frame}{تعريف المفهوم}
% TODO: اكتب التعريف الرسمي
\begin{block}{تعريف}
المفهوم هو وصف مختصر للتعريف الرسمي يوضح المعنى الأساسي
والمصطلح المقابل بالإنجليزية \textEN{(term)}.
\end{block}

\begin{itemize}
  \item السياق الذي يُستخدم فيه المفهوم.
  \item العناصر المكوّنة له.
\end{itemize}
\end{frame}

\begin{frame}{خصائص المفهوم}
% TODO: اذكر الخصائص الرئيسية
\begin{itemize}
  \item الخاصية الأولى وأثرها.
  \item الخاصية الثانية مع مثال موجز.
  \item الخاصية الثالثة \textEN{(property)}.
\end{itemize}
\end{frame}

% ============================================================
% SECTION 2 — EXAMPLE
% ============================================================
\section{مثال محلول}

\begin{frame}{مثال محلول}
% TODO: استبدل بمسألة فعلية من المقرر
\begin{block}{المسألة}
أوجد قيمة $x$ التي تحقق المعادلة التالية:
\[
2x + 5 = 15
\]
\end{block}

\textbf{الحل:}
\begin{align*}
2x + 5 &= 15 \\
2x &= 15 - 5 \\
2x &= 10 \\
x &= 5
\end{align*}
\end{frame}

% ============================================================
% SECTION 3 — APPLICATIONS
% ============================================================
\section{التطبيقات}

\begin{frame}{التطبيقات العملية}
% TODO: اذكر تطبيقات المفهوم
\begin{itemize}
  \item التطبيق الأول في المجال العملي.
  \item التطبيق الثاني \textEN{(use case)}.
  \item التطبيق الثالث مع قيود الاستخدام.
\end{itemize}
\end{frame}

% ============================================================
% SUMMARY
% ============================================================
\section{الملخص}

\begin{frame}{ملخص المحاضرة}
% TODO: لخّص النقاط الرئيسية
\begin{itemize}
  \item تعريف المفهوم وخصائصه الرئيسية.
  \item خطوات حل المسألة النموذجية.
  \item التطبيقات العملية وحدودها.
\end{itemize}
\end{frame}

% ============================================================
% FURTHER READING
% ============================================================
\begin{frame}{قراءات إضافية}
% TODO: أضف المراجع المقترحة
\begin{itemize}
  \item الفصل المرجعي من الكتاب المقرر.
  \item ورقة بحثية مرتبطة \textEN{(paper)}.
  \item مصدر إلكتروني للتعمق.
\end{itemize}
\end{frame}

\end{document}
